\section{A genuine action on the specified original Cartesian square}
\label{sec:pullback-cone-monoid-action}

\begin{theorem}[The actual Cartesian cone acquires the action {\cite[\S1.2]{BJ22}}]
\label{thm:pullback-cone-monoid-action}
\lean{EnvelopeProperty.PullbackConeMonoidAction.action, EnvelopeProperty.PullbackConeMonoidAction.action_fst, EnvelopeProperty.PullbackConeMonoidAction.action_snd, EnvelopeProperty.PullbackConeMonoidAction.modObj, EnvelopeProperty.PullbackConeMonoidAction.fst_isModHom, EnvelopeProperty.PullbackConeMonoidAction.snd_isModHom, EnvelopeProperty.PullbackConeMonoidAction.isModHom_of_comp, EnvelopeProperty.PullbackConeMonoidAction.lift_isModHom}
\leanok
Let $G$ be a monoid object in a Cartesian monoidal category.
Suppose the original diagram $X\xrightarrow f S\xleftarrow gY$
has genuine $G$-actions and equivariant maps.
Let $P\xrightarrow pX$, $P\xrightarrow qY$ be any specified
original Cartesian cone for this diagram. Then the \emph{same}
$P$ has an internally derived genuine $G$-action and both
original projections $p,q$ are equivariant. Equivariance of a
map into $P$ is detected by these two projections; in particular,
the original universal lift of two equivariant maps is equivariant.
\end{theorem}

\begin{proof}
Use the original Cartesian universal property to define
$\gamma_P:G\times P\to P$ by its two whole components
\[
 \gamma_Pp=(1_G\times p)\gamma_X,
 \qquad
 \gamma_Pq=(1_G\times q)\gamma_Y.
\]
Compatibility follows from equivariance of the original $f,g$,
functoriality of the product, and the original equality $pf=qg$.
To prove the unit law, postcompose both candidate maps with $p$
and $q$. The two component equations, naturality of the product
unit, and the original unit laws on $X,Y$ identify the resulting
maps. The original Cartesian uniqueness gives the whole unit law.
For the multiplication law use the same two projections,
naturality of the product associator, functoriality of products,
and the original multiplication laws on $X,Y$. This proves the
whole multiplication equation on $P$.

The defining component equations give equivariance of $p,q$.
For any original $a:T\to P$, if $ap$ and $aq$ are equivariant,
postcomposing the proposed equivariance equation for $a$ with
each original projection reduces it to those two equations.
Cartesian uniqueness then proves the entire equivariance equation.
The universal lift of two equivariant maps has exactly these
original two components, so the last assertion follows as well.
\end{proof}
